% Part: proof-theory
% Chapter: sequent-calculus
% Section: interpretation-rules

\documentclass[../../../include/open-logic-section]{subfiles}

\begin{document}

\olfileid[id]{pt}{seq}{irl}

\olsection{Interpretasi Aturan}

Ingat kembali interpretasi sekuen: sekuen $\Gamma\Sequent\Delta$ benar
dalam !!a{structure}~$\Struct{M}$ jika sekurang-kurangnya satu
$!A\in\Gamma$ salah atau sekurang-kurangnya satu $!A\in\Delta$ benar. Kita
dapat membaca aturan-aturan logis \Log{G3c} sebagai pengodean syarat
kebenaran !!{formula}s utamanya. Tinjau \RightR{\lif}. !!^{formula} utamanya
adalah $!A\lif !B$. Formula itu benar jika $!A$ salah atau $!B$ benar. Maka
sekuen $\ \Sequent !A\lif !B$ benar jika dan hanya jika
$!A\Sequent !B$ benar. Jika kita menambahkan !!{formula}s konteks
$\Gamma$, $\Delta$ masing-masing pada anteseden dan sukseden sekuen-sekuen
tersebut, kita memperoleh tepat kesimpulan dan premis aturan
\RightR{\lif}. Situasinya serupa bagi~\LeftR{\lif}. Formula
$!A\lif !B$ salah jika $!A$ benar dan $!B$ salah. Jadi,
$!A\lif !B\Sequent\ $ benar jika dan hanya jika baik $\ \Sequent !A$
maupun $!B\Sequent\ $ benar. Kesimpulan \LeftR{\lif} ekuivalen dengan
konjungsi premis-premisnya. Pola ini memastikan bahwa aturan \Log{G3c}
sahih (jika semua premis benar, kesimpulannya benar). Pola itu juga
memastikan kelengkapan.

Bahkan, syarat kebenaran setiap perangkai fungsional-kebenaran dapat
dideskripsikan dengan cara ini sebagai himpunan sekuen. Hal ini mengikuti
fakta bahwa setiap fungsi kebenaran dalam logika klasik dapat dideskripsikan
oleh !!a{formula} dalam bentuk normal konjungtif: konjungsi dari disjungsi
!!{formula}s atomik dan negasi formula atomik. Misalnya, tetapkan
$!A\oplus !B$ benar jika dan hanya jika tepat salah satu dari $!A$ dan $!B$
benar, yakni ``atau eksklusif''. Maka $!A\oplus !B$ benar jika dan hanya jika
$(!A\lor !B)\land(\lnot !A\lor\lnot !B)$ benar, dan formula itu salah jika
$(\lnot !A\lor !B)\land(!A\lor\lnot !B)$ benar. Karena itu, kita dapat
memperkenalkan aturan kalkulus sekuen bagi $\oplus$ sebagai
\begin{align*}
&
\Axiom$\Gamma \fCenter \Delta, !A, !B$
\Axiom$!A, !B, \Gamma \fCenter \Delta$
\RightLabel{\RightR{\oplus}}
\BinaryInf$\Gamma \fCenter \Delta, !A \oplus !B$
\DisplayProof
\\
& \Axiom$!A, \Gamma \fCenter \Delta, !B$
\Axiom$!B, \Gamma \fCenter \Delta, !A$
\RightLabel{\LeftR{\oplus}}
\BinaryInf$!A \oplus !B, \Gamma \fCenter \Delta$
\DisplayProof
\end{align*}
dan aturan-aturan ini juga sahih dan lengkap. Kesimpulan setiap aturan benar
jika dan hanya jika semua premisnya benar.

\begin{prob}
Andaikan $!A\downarrow !B$ benar jika dan hanya jika baik $!A$ maupun $!B$
tidak benar (``NOR''). Bagaimanakah aturan \Log{G3c} bagi perangkai tersebut?
(Temukan dua sekuen yang benar secara bersamaan jika dan hanya jika
$!A\downarrow !B$ benar. Hal itu memberikan aturan \RightR{\downarrow}.
Kemudian temukan satu sekuen yang benar jika dan hanya jika
$!A\downarrow !B$ salah. Hal itu memberikan aturan \LeftR{\downarrow}.)
\end{prob}

Dalam \Log{G3c}, setiap perangkai dan kuantor mempunyai tepat dua aturan:
aturan kiri dan aturan kanan. Namun, dalam \Log{G1c} terdapat dua aturan
\LeftR{\land} dan dua aturan \RightR{\lor}. Alasannya, sistem asal Gentzen
\Log{LK}, yang menjadi model bagi \Log{G1c}, mempunyai varian~\Log{LJ} bagi
suatu logika nonklasik yang penting, yaitu logika intuisionistik. Untuk
mendapatkan sistem yang sahih dan lengkap bagi logika intuisionistik,
sukseden setiap sekuen dalam~\Log{LJ} dibatasi agar memuat paling banyak satu
!!{formula}. Hal ini membuat aturan \RightR{\lor} dari~\Log{G3c} tidak dapat
digunakan, karena premisnya memuat baik $!A$ maupun $!B$ dalam sukseden.
Karena itu, Gentzen menggunakan sepasang aturan \RightR{\lor} bagi
\Log{LJ}, dan menggunakan aturan yang sama bagi~\Log{LK}. Versi
intuisionistik~\Log{G1i} juga hanya merupakan~\Log{G1c} dengan semua sukseden
dibatasi agar memuat paling banyak satu !!{formula}. (\Log{G3c} juga
mempunyai versi intuisionistik, tetapi beberapa aturannya berbeda dari aturan
yang bersesuaian dalam~\Log{G3c}. Misalnya, versi itu juga menggunakan
sepasang aturan \RightR{\lor}.)

Bahwa aturan-aturan struktural \Log{G1c} sahih mudah dilihat. Misalnya, jika
$\Gamma\Sequent\Delta$ memuat !!{formula} salah di sebelah kiri atau
!!{formula} benar di sebelah kanan, maka $\Gamma\Sequent\Delta,!A$ juga
memuat formula yang sama; benar atau salahnya $!A$ tidak relevan. Jadi, jika
premis \RightR{\Weakening} benar, kesimpulannya juga benar. Perhatikan bahwa
konversnya tidak benar. Kesimpulan mungkin benar karena $!A$ benar, sedangkan
premisnya tidak dijamin benar.

Mengapa aturan struktural diperlukan? Kontraksi
(\RightR{\Contraction}, \LeftR{\Contraction}), misalnya, diperlukan untuk
!!{prove} $\ \Sequent !A\lor\lnot !A$. Jika kita mulai dengan
$!A\Sequent !A$, mula-mula kita memperoleh $\Sequent !A,\lnot !A$ melalui
\RightR{\lnot}. Kemudian kita dapat menggunakan dua versi
\RightR{\lor} secara berurutan: sekali untuk memperoleh
$!A\lor\lnot !A$ dari $\lnot !A$, dan sekali untuk memperoleh
$!A\lor\lnot !A$ dari~$!A$. Hasilnya adalah
$\ \Sequent !A\lor\lnot !A,!A\lor\lnot !A$. Jadi, untuk memperoleh
$\Sequent !A\lor\lnot !A$ dalam~\Log{G1c}, kita harus dapat
``mengontraksikan'' dua kemunculan !!{formula} yang sama dalam !!a{proof}:
\[
\Axiom$!A \fCenter !A$
\RightLabel{\RightR{\lnot}}
\UnaryInf$\fCenter !A, \lnot !A$
\RightLabel{\RightR{\lor}}
\UnaryInf$\fCenter !A, !A \lor \lnot !A$
\RightLabel{\RightR{\lor}}
\UnaryInf$\fCenter !A \lor \lnot !A, !A \lor \lnot !A$
\RightLabel{\RightR{\Contraction}}
\UnaryInf$\fCenter !A \lor \lnot !A$
\DisplayProof
\]
Memang terdapat sekuen valid yang tidak dapat kita !!{prove} dalam~\Log{G1c}
tanpa pelemahan atau kontraksi. Misalnya, !!a{proof} bagi sekuen
$!A\Sequent !B\lif !A$ dalam~\Log{G1c} harus berakhir dengan
\RightR{\lif} (karena hanya $!B\lif !A$ yang dapat menjadi formula utama),
dan aturan itu memerlukan $!B,!A\Sequent !A$ sebagai premis. Premis tersebut
pada gilirannya memerlukan \LeftR{\Weakening} agar dapat diperoleh dari
aksioma $!A\Sequent !A$:
\[
\Axiom$!A \fCenter !A$
\RightLabel{\LeftR{\Weakening}}
\UnaryInf$!B, !A \fCenter !A$
\RightLabel{\RightR{\lif}}
\UnaryInf$!A \fCenter !B \lif !A$
\DisplayProof
\]

Dalam \Log{G3c}, aturan kontraksi dan pelemahan sama sekali tidak diperlukan.
Pelemahan sudah tercakup dalam aksioma. Kontraksi sebagian besar dapat
dihindari dengan menggabungkan dua salinan konteks $\Gamma$, $\Delta$ dalam
aturan berpremis dua menjadi satu. Baik \Log{G1c} maupun \Log{G3c} mempunyai
aturan ``berbagi konteks'' semacam itu. Dalam aturan berpremis satu, kontraksi
hanya diperlukan bagi \RightR{\lor}, \LeftR{\land}, \RightR{\lexists},
dan~\LeftR{\lforall}. Versi aturan-aturan tersebut dalam~\Log{G3c} membuat
kontraksi sepenuhnya berlebihan. Terdapat sistem \Log{G2c} (lihat
\olref{tab:G2c}) yang aturan-aturan berpremis duanya mempunyai konteks
independen. Dalam \Log{G2c}, kontraksi bahkan lebih penting daripada
dalam~\Log{G1c}.

\subfile{rules-G2c}

\end{document}
